% =============================================================================
%  Upper Volume, Articles 16--42
%  Generated from translation/pages_31-34.md, pages_35-38.md, pages_39-42.md
% =============================================================================

% -----------------------------------------------------------------------------
%  Article 16
% -----------------------------------------------------------------------------
\article{16}{The Range-Bound Market after a Ceiling Price}

\begin{sokyubox}
{\small\itshape\color{sokyubar} \jp{第十六章\quad 天井直段後通相塲の事}}

\medskip
After a ceiling price (\jt{天井直段}{tenj\=o nedan}), in a range-bound
market (\jt{通相塲}{kayoi s\=oba}) that is trending downward, prices
rise at the opening of the month and decline by the twenty-ninth
or thirtieth, falling by twenty to thirty or thirty to forty
\jt{俵}{tawara} without fail. Of this there is no doubt.
\end{sokyubox}

\commentary
In a falling market after a ceiling, prices in the first ten days
of the month tend to edge upward, because the memory of the
previous month's high price lingers and there is much wistful
buying (\jt{買慕ひ}{kaimurahi}). However, this is all it amounts to.
The surrounding conditions offer no bullish factors; rather, with
arrivals of physical rice (\jt{正米}{sh\=omai}), with producing regions
reporting low prices, with favorable weather, and with other
bearish factors at work, buyers naturally thin out and selling
orders (\jt{投物}{nagemono}) appear. By mid-month or the latter part
of the month, prices invariably decline. One should therefore
consult the chart (\jt{罫線表}{keisen-hy\=o}), confirm that the
post-ceiling range-bound rally has occurred, and establish a sell
strategy.

% -----------------------------------------------------------------------------
%  Article 17
% -----------------------------------------------------------------------------
\article{17}{A Seventh- or Eighth-Month Ceiling Followed by Five or Six Months of Decline}

\begin{sokyubox}
{\small\itshape\color{sokyubar} \jp{第十七章\quad 七八月天井五六ケ月引下の事}}

\medskip
When a ceiling price appears in the seventh or eighth month, prices
decline through the twelfth month. In the first month they rise
slightly, but through the spring the market remains sluggish. If,
during this period, prices have fallen excessively, they will rise
sharply at some point during the fifth or sixth month.
\end{sokyubox}

\commentary
A ceiling price around the seventh or eighth month is most often
caused by wind, rain, flooding, and the like---the high price
having been driven to ceiling levels by such events. After the
fact, when the weather returns to normal and crop prospects improve,
prices decline gradually, continuing to fall for as long as five or
six months through the twelfth month. Around the first month,
profit-taking (\jt{利喰}{rigui}) purchases emerge, and with prices
considerably below the ceiling, speculative buying also appears, so
prices rise somewhat. Nevertheless, the rice is the tail end of a
bumper crop; buyers recall the previous year's experience and are
reluctant to press their purchases aggressively. The market
therefore fails to rally through the spring and consolidates
(\jt{保合}{mochai}) at low levels. However, by the fourth or fifth
month, complaints begin to surface about oversupply from producing
regions or a poor wheat harvest, and the market creeps upward,
sometimes continuing to rise through the fifth or sixth month.

% -----------------------------------------------------------------------------
%  Article 18
% -----------------------------------------------------------------------------
\article{18}{Rice Cheap through the Eleventh, Twelfth, and First Months Rallies in Summer}

\begin{sokyubox}
{\small\itshape\color{sokyubar} \jp{第十八章\quad 十一十二正月頃迄下直夏引上の事}}

\medskip
Rice that remains cheap through the eleventh month, twelfth month,
and first month should be understood as destined for a summer rally.
It may rise through the seventh month. Take note of this.
\end{sokyubox}

\commentary
The previous chapter explained why rice that made its ceiling in
the seventh and eighth months declines through the eleventh and
twelfth months. When the lingering momentum of the previous year's
bumper harvest, together with pressure from physical rice
deliveries and the like, pushes prices down from the eleventh month
through the first month, such a market has already voiced every
bearish argument there is to voice and has fallen as far as it can
fall. Its rebound is therefore inevitable. In particular, after a
period in which bullish factors have been exhausted through
over-praising, conditions tend to turn bullish again; and when some
fault in the spring wheat crop is reported, the rally can
sometimes extend from the fifth month through the entirety of the
sixth month.

% -----------------------------------------------------------------------------
%  Article 19
% -----------------------------------------------------------------------------
\article{19}{A Seventh- or Eighth-Month Floor Leading to a Major Rally through the Twelfth or First Month}

\begin{sokyubox}
{\small\itshape\color{sokyubar} \jp{第十九章\quad 七八兩月応十二正月迄大上けの事}}

\medskip
Rice at its floor price (\jt{底直段}{soko nedan}) in the seventh and
eighth months should produce a major rally by the twelfth or first
month.
\end{sokyubox}

\commentary
Rice at its floor in the seventh and eighth months is, in the
final analysis, the result of the previous year's good harvest
being reinforced by a strong wheat crop, which together drove
prices down; then, with favorable weather being praised before and
after the midsummer season (\textit{doy\=o}), and crop conditions appearing
exceptionally good, the market foresaw a bumper harvest and voiced
every bearish argument until prices were pressed to their lowest.
However, as a matter of market rhythm, the rebound must come;
meanwhile, natural disasters begin to occur, and reports of
insufficient sickle entry (\jt{鎌入不足}{kamairi-fusoku}---an
inadequate harvest yield) begin to circulate. The more the crop
was over-praised at the outset, the more violent the rebound. This
rally can sometimes continue through the first month.

% -----------------------------------------------------------------------------
%  Article 20
% -----------------------------------------------------------------------------
\article{20}{The Months in Which Prices Rise from the Floor}

\begin{sokyubox}
{\small\itshape\color{sokyubar} \jp{第二十章\quad 底より起き上り月の事}}

\medskip
The ninth, tenth, and eleventh months---these three months
are emphatically months in which no ceiling price appears. They
are months in which prices rise from the floor.

However, notwithstanding the above, there are years when the
ceiling appears in the eighth, ninth, or tenth month. Consider
this carefully.
\end{sokyubox}

\commentary
According to the approximate methods for estimating ceiling and
floor months described in the preceding chapters, the three months
of the ninth, tenth, and eleventh fall right in the middle of the
cycle. This is why the Master stated the principle broadly in this
way. However, the examples given in the previous chapters all
concern six-month spans. If the span shortens to three months or
more but under five months, then these same three months could
correspond to ceiling months, so one cannot state the principle as
absolutely as this chapter does. In essence, though, during these
three months the disaster season has passed and the approximate
shape of the harvest is known. If conditions prove uneventful, it
is only natural that the floor price should appear in these three
months. This is presumably why the Master stated the principle in
such broad terms.

% -----------------------------------------------------------------------------
%  Article 21
% -----------------------------------------------------------------------------
\article{21}{How to Read a Rising Market}

\begin{sokyubox}
{\small\itshape\color{sokyubar} \jp{第二十一章\quad 上相塲見様の事}}

\medskip
The timing of the ceiling in a rising market is not fixed. In
general, the month's floor appears early in the month, and
the month's ceiling emerges between the twenty-first and the
twenty-sixth. Rice that is rising from the floor continues to
rise for many months. When it has gradually risen by a hundred
bales or more, and then at month-end---through the twenty-ninth
and thirtieth---it still edges up by ten, twenty, or even thirty
bales, know that the following month will produce the year's
annual ceiling price (\jt{年中行付天井直段}{nenj\=u yukitsuki tenj\=o
nedan}). Of this there is no doubt.
\end{sokyubox}

\commentary
The method for observing the timing of ceilings and floors in a
rising market has no fixed rule. However, the principle stated
repeatedly in earlier chapters---that in a rising market, the
month opens weak and closes strong---is beyond doubt. In general,
the month's floor price appears early in the month, and between
around the twentieth and the twenty-sixth or twenty-seventh, prices
surpass the previous month's monthly high-low mark (\jt{月節}{tsukibushi})---that is, the previous high---to set the new
monthly high. Moreover, rice that is rising from the floor
oscillates in this pattern---advancing, retracing, advancing,
retracing---and continues to rise over several months. When,
over the course of these several months, prices have risen by a
hundred bales, and then on top of that, at month-end, they climb
another ten, twenty, or even thirty bales in a hissing rush
(\jp{ヒス〳〵}, \textit{hisuhisu}), the following month will without fail
produce the year's highest price.

% -----------------------------------------------------------------------------
%  Article 22
% -----------------------------------------------------------------------------
\article{22}{A Ceiling Consolidation from the First Month through the Third or Fourth Month}

\begin{sokyubox}
{\small\itshape\color{sokyubar} \jp{第二十二章\quad 正月より三四月迄天井持合の事}}

\medskip
Rice that has been consolidating at ceiling-price levels from
around the first month through the third or fourth month will
assuredly break lower during the fifth or sixth month. If it
declines fully in the fifth month, the sixth month will bring a
sharp rally. If it does not decline in the fifth month, know that
it will assuredly decline in the sixth month. Of this there is no
doubt. A fifth- or sixth-month decline should be closed out
cleanly (\textit{sappari}) by the twenty-first. Declines in other months
run through the twenty-ninth or thirtieth.
\end{sokyubox}

\commentary
This chapter explains the summer break that follows a spring
consolidation (\jt{保合}{mochai}).

Rice that produced a ceiling price by the twelfth or first month
due to a poor harvest will, even after physical rice from supply
regions has been drawn out in large volume by the high futures
price, retain bullish sentiment (\jt{人氣}{ninki})---people remember
the previous year's poor harvest and buying interest remains
strong; every dip is targeted for purchase by the crowd. A major
decline therefore does not come easily. Nevertheless, in such
cases, one must not forget the sell strategy premised on a
post-poor-harvest ceiling consolidation. The reason is that this
consolidation never lasts indefinitely. Eventually the fifth month
arrives; the wheat crop proves average; the weather turns
favorable; and on top of that, the rice from the previous poor-harvest
year is of inferior quality, so there is anxiety about
spoilage in the warm weather, and rice holders begin to sell
urgently. The result is a sudden, severe collapse.

However, the rule of thumb for the fifth- and sixth-month collapse
is this: if one has profits, take them all by mid-month---around
the twenty-first. The reason is that in other months this need not
apply, but as the fifth month draws to a close, the new trading
period begins in the sixth month, and seller profit-taking will
emerge. Moreover, the sixth month gradually enters the natural-disaster
season, making the sell side hazardous.

% -----------------------------------------------------------------------------
%  Article 23
% -----------------------------------------------------------------------------
\article{23}{Price Swings in the Fifth and Sixth Months}

\begin{sokyubox}
{\small\itshape\color{sokyubar} \jp{第二十三章\quad 五六月兩月相場高下の事}}

\medskip
From mid-fifth month through the sixth month, when rice rises
sharply, it will also fall sharply. When it falls sharply, the
same holds for the rally. The same applies to new-crop rice during
the seventh, eighth, and ninth months.
\end{sokyubox}

\commentary
This chapter illustrates, by way of example, that market sentiment
(\jt{人氣}{ninki}) tends to overreact to changing conditions and
panic in fear.

Now, from around mid-fifth month through the ninth month lies the
natural-disaster season. When word comes of flooding, people rush
to buy without properly estimating the extent of the damage---a
fear-driven, sudden rally. Then, when conditions appear to improve,
a reactionary sell-off follows, producing a chaotic, disordered
market. The same holds true of typhoons. In short, the market
during the disaster season is one in which a sharp rally is
followed by a sharp decline, and a sharp decline is followed by a
sharp rally. However, this is a broad generalization that holds in
eight or nine cases out of ten.

% -----------------------------------------------------------------------------
%  Article 24
% -----------------------------------------------------------------------------
\article{24}{The Sell-Off of Rice in a Good-Harvest Year}

\begin{sokyubox}
{\small\itshape\color{sokyubar} \jp{第二十四章\quad 上作年の米売崩しの事}}

\medskip
In a year when both \jp{最上} (Mogami) and our own region
have a good harvest, if rice has been holding at a relatively high
price through the winter and into the first month, a sell-off
should occur between mid-third month and mid-fourth month.
\end{sokyubox}

\commentary
``Mogami'' refers to the neighboring province of Sh\=onai---the
present-day Murayama and Mogami districts. The meaning is that
even in a year when both nearby provinces and the local region have
had a good harvest, if for some reason the price has been holding
at a relatively high level compared to other markets, eventually
physical rice will arrive, and once prices dip slightly, a wave of
speculative selling and new short positions will appear all at
once. The result is a major collapse around the time of the
equinox (\jt{彼岸}{higan}).

% -----------------------------------------------------------------------------
%  Article 25
% -----------------------------------------------------------------------------
\article{25}{Maneuvering Ahead of Crop-Condition Reports}

\begin{sokyubox}
{\small\itshape\color{sokyubar} \jp{第二十五章\quad 作合取沙汰前方掛引の事}}

\medskip
Reports of crop conditions and complaints begin to emerge around
mid-seventh month. Even when a good harvest has been proclaimed,
some manner of adverse report will inevitably surface. Maneuvering
(\jt{掛引}{kakehiki}) in advance of this is of the utmost importance.

``Complaints'' (\jt{申分}{m\=oshibun}) means crop damage or faults.
\end{sokyubox}

\commentary
The meaning is this: even when favorable weather has been praised
and the crop has been extolled, from around mid-seventh month
onward, reports begin to emerge---flooding in one area, insect
damage in another, a low-pressure system forming in yet another.
The market tends to lean toward the upside. Therefore, anyone who
has sold short in advance should be prepared, and it is essential
to maneuver (\textit{kakehiki}) toward closing the position.

% -----------------------------------------------------------------------------
%  Article 26
% -----------------------------------------------------------------------------
\article{26}{Carrying a First-Month Profit into the Second Month}

\begin{sokyubox}
{\small\itshape\color{sokyubar} \jp{第二十六章\quad 正月利運持越しの事}}

\medskip
A profitable rice position from the first month, if carried over
into the second month, will prove most unsatisfactory. One should
close it out cleanly within the first month. However, a position
established in the first month based on one's reading of the
twelfth month's conditions is a separate matter.
\end{sokyubox}

\commentary
As the following chapter explains, the market in the first and
second months tends toward a simple back-and-forth: if the first
month rises, the second month falls; if the first month falls,
the second month rises. Therefore, if one holds a profitable
position from the first month, one should close it out promptly.
This, however, is stated with reference to the market's normal
monthly rhythm; it is not impossible for the first and second
months to move in the same direction. But that can be discerned
from the market's behavior in the twelfth month, and a trade
(\jt{商内}{akinai}) founded on that reading is excepted from this
rule.

% -----------------------------------------------------------------------------
%  Article 27
% -----------------------------------------------------------------------------
\article{27}{The Dull First and Second Months, the Fourth-Fifth-Sixth Month Collapse}

\begin{sokyubox}
{\small\itshape\color{sokyubar} \jp{第二十七章\quad 正二売買退屈四五六崩しの事}}

\medskip
The market in the first and second months sees no great swings.
However, the pattern for these two months is: if the first
month rises, the second month falls; if the first month falls, the
second month rises. One should regard each as a single swing of
the pendulum.
\end{sokyubox}

\commentary
The market in the first and second months sees no great swings, and
is mostly a range-bound consolidation (\jt{持合相場}{mochai s\=oba}).
When prices move even slightly in either direction, people react
as though the signal has arrived, and sentiment follows in a rush
of blind imitation. Yet because there is no real substance behind
the move, the advance quickly loses momentum. The observation
that if the first month rises the second month falls, and if the
first month falls the second month rises, amounts to a description
of this effect of sentiment. However, this is a general principle
and not a fixed rule; hence the emphatic warning: ``Above all, one
must watch for the latent market.'' The latent market
(\jt{潜相場}{hisomi s\=oba}) refers to a
consolidation that is on the verge of breaking in one direction
or the other---a situation in which, as the saying goes, ``the
mountain rain is about to come and the wind fills the tower,'' where
the atmosphere is charged with the potential for a sudden move.
To perceive this, one should consider whether the current price
stands at the ceiling, or at the floor, or whether it is a
consolidation at a middle level. In short, speaking from the
normal tendency of the market, the first and second months are
months of small swings, and both buyers and sellers must steel
their nerves. The third month is a strong month. The fourth, fifth,
and sixth months are months when the market tends to collapse.

% -----------------------------------------------------------------------------
%  Article 28
% -----------------------------------------------------------------------------
\article{28}{Floor and Ceiling of the November Contract}

\begin{sokyubox}
{\small\itshape\color{sokyubar} \jp{第二十八章\quad 霜月限底天井の事}}

\medskip
When a new November contract (\jt{霜月限}{shimotsuki-giri}) opens and
rice begins to rise from a floor of around forty-five or forty-six
to fifty bales, if even a minor crop complaint (\jt{申分}{m\=oshibun})
emerges, the ceiling price should not be regarded as limited to a
hundred-bale rise. Depending on circumstances, one should expect
rises of a hundred sixty or even a hundred eighty bales. However,
when rice begins to rise in the sixth or seventh month from a level
of only thirty-four or thirty-five bales---that is, from a
mid-level floor---then even if the local crop conditions are
favorable and even if carry rice (\jt{乗米}{nosei-mai}) is flowing in
to lift the Osaka market, a hundred-bale rise should be regarded
as the ceiling.
\end{sokyubox}

\commentary
Rice that begins to rise from a major floor of around fifty bales,
when even a slight crop deficiency appears, will not stop at a
hundred-bale rise but may reach a ceiling of a hundred sixty or
even a hundred eighty bales above. But rice that starts upward
from a mid-level floor need only rise a hundred bales before one
should regard it as having reached the ceiling.

This chapter explains that a market rising from a middle-level
position will produce a smaller total advance to the ceiling,
whereas a market rising from a true floor will produce
proportionally larger gains. To give an example: if the advance
from a middle level is fifteen percent, then from a major floor
the advance may reach as much as thirty percent.

% -----------------------------------------------------------------------------
%  Article 29
% -----------------------------------------------------------------------------
\article{29}{On the Range-Bound Market After a Major Swing}

\begin{sokyubox}
{\small\itshape\color{sokyubar} \jp{第二十九章\quad 大高下過通ひ相塲の事}}

\medskip
After a major swing has run its course and the ceiling price has
been struck, the market enters a consolidation (\jt{保合}{mochai}),
waiting to see whether it will go up or down. When, amid this
uncertainty, the market takes on a vaguely bullish tinge---perhaps
on reports from the Kamigata region---and sentiment brightens a
little, everyone turns buyer and the market rises enough to cause
excitement. At that point, it is decidedly a selling opportunity.
This is a range-bound market (\jt{通相場}{kayoi s\=oba}): it rises and
then falls, falls and then rises, without any excessive swing, and
oscillates back and forth repeatedly. One must not let one's heart
be swayed by sentiment or by the oscillations of the range-bound
market; considering the ceiling and the floor is of first
importance. However, when rice consolidates at the floor without
any great swing and begins to rise naturally, one must not sell it.
Rather, one should add to one's buying position step by step. This
is not a range-bound market---it is a market that is rising of its
own accord.
\end{sokyubox}

\commentary
This chapter sets out the principles for selling into the
range-bound market that follows a ceiling. The meaning is
sufficiently clear without further annotation.

After the ceiling, range-bound rallies may present selling
opportunities for several months, but one should estimate the
elapsed months and days since the ceiling and floor, together with
the magnitude of the price range, and avoid overcommitting. In
particular, a market that has established its floor and begins to
rise naturally can sometimes be mistaken for a range-bound market.
But one must never sell into a market that rises naturally without
great swings; on the contrary, one should add to one's buying
position.

% -----------------------------------------------------------------------------
%  Article 30
% -----------------------------------------------------------------------------
\article{30}{Secret Tradition of Buying and Selling After a Consolidation}

\begin{sokyubox}
{\small\itshape\color{sokyubar} \jp{第三十章\quad 持合後売買秘傳の事}}

\medskip
During the above-described consolidation, when the market turns
slightly downward, everyone concludes that this is the beginning
of a bear market. Those already short congratulate themselves and
press their selling further. Even the longs abandon ship and sell
to escape, then find themselves overselling. Everyone rushes to sell
first, and the market falls still further. At this moment, one
should buy. The profit is assured. It is an extremely difficult
purchase to make, yet buy one must. Even experienced traders of
many years often regret missing it. In range-bound trading, one
should target a ten-bale swing and take profits quickly---that is
the first principle.
\end{sokyubox}

\commentary
In the consolidation following a ceiling, as described in the
preceding chapter, when the market turns even slightly downward,
the general opinion shifts to bearishness. The shorts, feeling that
the market has come over to their side, ride their advantage and
hammer the price down---and the longs, too, scramble to sell and
escape, falling over one another to liquidate. The market
accordingly falls further. But at precisely such a moment, one
should buy. It is an extremely difficult purchase---even seasoned
traders find it hard to pull the trigger---yet if one buys, the
profit is certain. However, the typical range-bound swing amounts
to roughly ten bales, so one should estimate accordingly and take
profits nimbly.

% -----------------------------------------------------------------------------
%  Article 31
% -----------------------------------------------------------------------------
\article{31}{On the Discipline When a Buy Opportunity Is Missed}

\begin{sokyubox}
{\small\itshape\color{sokyubar} \jp{第三十一章\quad 可買入見合候時心得の事}}

\medskip
When one has identified rice as a buy, but the price has already
risen two bales or so and one feels one has been left behind, one
sometimes reverses course and turns seller out of frustration. This
is a grave error. When one has missed the entry, one should simply
wait for the next buying opportunity.
\end{sokyubox}

\commentary
In forming a market view, there are two faculties one must
understand: one is intellect (\jt{智}{chi}), and the other is emotion
(\jt{情}{j\=o}). A forward-looking assessment should be decided entirely
by intellect; one must never, ever allow emotion to intrude.

So too in the case of this chapter: to identify rice as a buy was
an act of intellect. But to hesitate, then to watch the price rise
two bales, then to regret the missed opportunity and rashly turn
seller---this is to discard the intellect that is essential to
sound judgment and to adopt the emotion that invariably leads
judgment astray. It is as though one were to throw away jewels and
pick up dung. Having formed the view that one should buy, one
should hold to that original assessment and wait for the
range-bound pullback---that is, for the dip.

% -----------------------------------------------------------------------------
%  Article 32
% -----------------------------------------------------------------------------
\article{32}{When the Market Has Been Pulled Up Relentlessly for Two or Three Months}

\begin{sokyubox}
{\small\itshape\color{sokyubar} \jp{第三十二章\quad 二三ケ月必至々々と引上候時の事}}

\medskip
When the market has been rising relentlessly for two or three
months and then suddenly drops by about one fall (\jt{一落}{hitootoshi})---that is, three or four bales---one should
decidedly buy. It will rise again. At that point, one should
promptly take profits. This is maneuvering (\jt{駈引}{kakehiki}) of
the first order, and it is extremely difficult to execute.
\end{sokyubox}

\commentary
When the market has been rising vigorously for two or three months
and then suddenly plunges three or four bales (one fall), such a
sudden drop within an ongoing advance is a range-bound correction;
it will necessarily retrace to some extent, so one should buy. But
whether the advance will end after two or three months, or whether
the market will continue to oscillate and keep climbing for another
four or five months, cannot be known until one sees what follows.
Therefore one should take profits promptly and reassess. However,
when one has a small profit in hand, human weakness is to ride one's
luck; this kind of maneuvering is consequently very hard to carry
out.

% -----------------------------------------------------------------------------
%  Article 33
% -----------------------------------------------------------------------------
\article{33}{On Closing a Trade and Resting Forty or Fifty Days}

\begin{sokyubox}
{\small\itshape\color{sokyubar} \jp{第三十三章\quad 商仕舞四五十日休の事}}

\medskip
When one has bought rice on the conviction that it will rise, and
the market advances beyond one's original expectation and one's
trade has paid off handsomely---then, at the point where the
advance stalls and runs into resistance, the market will swing
erratically. At that point, one should close the trade and rest for
forty or fifty days. During this rest, even if an opportunity
appears so compelling that one feels one can see through the
market, one must absolutely not trade. Regardless of the volume of
carry rice (\jt{乗米}{nosei-mai}) or the inflow of ship-borne
purchases, when a trade has paid off fully, one should close it out
at that juncture, rest for forty or fifty days, carefully assess
the strength or weakness of the rice market, compare it against the
Three-Position Tradition (\jt{三位の傳}{sanmi no den}), determine the
proper stance, and only then trade. The same applies when one has
profited on the sell side.
\end{sokyubox}

\commentary
This chapter admonishes against riding one's luck and chasing the
market after a winning trade.

Now, when one has bought into a rising market and one's forecast
has been entirely correct, and the advance stalls and runs into
resistance, the market will swing wildly. Because one's forecast
has been right, one feels reluctant to close the trade and wants
to ride it further. But at such a time, one should take profit on
every lot and rest for forty or fifty days. Observing the market
from the sidelines of a neutral position, one sees the swings with
far greater clarity. During that rest, however, one must absolutely
not trade. It does not matter if there is ample margin between
regional prices, or if buy orders are flooding in from elsewhere---none of these circumstances should concern one. When a trade has
fully paid off, one takes profit, rests for a month or two, calmly
assesses the market's condition and the direction of sentiment, and
with a clear and composed mind sets one's course before resuming
trading. The same applies when one has profited on the sell side.

% -----------------------------------------------------------------------------
%  Article 34
% -----------------------------------------------------------------------------
\article{34}{On the Discipline When Profits Accrue on a November Contract Purchase}

\begin{sokyubox}
{\small\itshape\color{sokyubar} \jp{第三十四章\quad 霜月限買入利分付候時心得の事}}

\medskip
In the November contract, when one has identified the floor, bought
in, and a considerable profit has accrued, the market sometimes
enters a consolidation or pulls back slightly. At that point, one
calculates the carrying cost (\jt{利足}{risoku}) and regrets not
having sold on the earlier rally. This is a grave
misunderstanding. When one has bought at the floor, one must never
sell before a margin call (\jt{落引}{ochibiki}) has occurred. From the
floor purchase through the advance, one should continue to add to
one's position until the margin call comes. Keep this well in mind.
\end{sokyubox}

\commentary
The ``carrying cost'' (\textit{risoku}) refers to a custom of the Sakata
Rice Exchange (\jt{酒田米會所}{Sakata kome kaisho}): when a trade is
not closed within the month in which it was opened and is carried
over to the following month, the buyer must pay to the seller, at
the time of profit-and-loss settlement, a sum reckoned at a fixed
rate per month. The details are set out in the section on the
customs of the Sakata Rice Exchange at the beginning of this
volume.

\begin{sokyubox}
{\small\itshape\color{sokyubar} Article 34 (continued)}

\medskip
This chapter discusses maneuvering on the buy side. The point is
this: when one has bought at the floor and the trade is in profit,
and the market dips slightly, one thinks: if only I had taken
profit when the market was higher the other day, I would not have
to pay the carrying cost for rolling the position into the next
month---and one feels regret. But this is a misunderstanding.
Rice that begins to rise from the floor will not reach the ceiling
until it has risen at least ten bales; therefore, until it has
risen at least one fall---that is, three or four bales---one must
not sell. On the contrary, one may add to the position.
\end{sokyubox}

% -----------------------------------------------------------------------------
%  Article 35
% -----------------------------------------------------------------------------
\article{35}{The Heart of the Lamp That Flares Before It Dies}

\begin{sokyubox}
{\small\itshape\color{sokyubar} \jp{第三十五章\quad 燈火消えんとして光増心の事}}

\medskip
In the November contract, when the price is near the ceiling and
one holds a long position, it sometimes happens that no significant
new commitments materialize, yet the market surges upward
regardless. Even a trader who has bought two or three hundred
\textit{ry\=o} worth regrets not having bought more and feels vexed; he
thinks he will buy the next small dip, but each day the market
climbs higher, margin calls begin to appear, and positions taken
just two or three days earlier are already being called. Thinking
that he may be late but can still capture at least half a fall's
profit, he adds to his position in rice that has already been
driven to the top. This is extremely ill-advised and must be
avoided. There is no profit to be taken from buying at this level.
Even if one buys a hundred \textit{ry\=o} worth, the profit from the
original cheap purchase of two hundred \textit{ry\=o} will be wiped out,
leaving not merely no gain but an actual loss. Once one buys at
such a level, one will be tempted to average down by buying more
on the next small dip, compounding the damage. Furthermore, even
if bullish reports of every kind arrive from the Kamigata region,
one must never buy into the commotion. Keep this well in mind.

This kind of rally climbs and climbs until it seems it might rise
forever. This is the heart of the lamp that flares just before it
goes out (\jp{燈火消えんとして光増す}, \textit{t\=oka kien to shite hikari
masu}). Even if one is utterly convinced that tomorrow will bring
further profit, one must not buy. At this point, the rice is
strong and sentiment is strong, and bullish conviction swells
beyond all measure. It is precisely here that one must watch for
the terminal ceiling (\jt{行付天井}{yukitsuki tenj\=o}). Consider this
carefully.
\end{sokyubox}

\commentary
This chapter warns against being swept into the vortex when the
price is near the ceiling and everyone---without exception---has
turned bullish.

Now, when the November contract approaches the ceiling and one's
position is in profit, it sometimes happens that no great volume of
new commitments is produced, yet the market rises sharply. Even a
trader holding a hundred lots regrets not having bought two or
three hundred and feels vexed; he waits to buy the next small dip,
but each day the market edges higher, and the positions he bought
just two or three days ago are already showing heavy profits.
Thinking that, though he is a little late, there is still something
to be gained at these levels, he adds to his position---unaware
that the market has already reached its terminal ceiling. But this
must be avoided. Buying at such levels never results in profit.
Even if one buys a hundred lots, the profit from the initial cheap
purchase of two hundred lots will be consumed by the high-priced
addition, and when the market eventually dips, the result is not
merely no profit but outright loss. Once one buys at such a level,
stubbornness sets in: on the next small dip, one wants to buy
still more to average down. And when bullish reports and all manner
of supporting reasons come in from the Kamigata region, sentiment
grows lively and excitable---but one must never buy into such
commotion. Bear this firmly in mind.

The kind of rally described in this chapter is a market that has
climbed to the extreme of its bullish phase---the turning point is
already at hand. Sentiment is such that the bulls pile on their
buying, and the bears, shaken and fearful, feel the market might
rise without limit. This is what the proverb means when it says,
\textit{``A lamp burns brightest just before it goes out.''} Even when it
seems an absolute certainty that tomorrow will bring still higher
prices, and even though buying here appears a sure profit as
obvious as picking up coins, one must never buy. At this juncture,
sentiment is strong and the market is strong, and bullish traders
are desperate to buy more---but it is precisely here that fixing
one's eye on the ceiling is the great secret of the market
practitioner.

% -----------------------------------------------------------------------------
%  Article 36
% -----------------------------------------------------------------------------
\article{36}{On a Falling Market within a Rising Trend}

\begin{sokyubox}
{\small\itshape\color{sokyubar} \jp{第三十六章\quad 上の内の下相場の事}}

\medskip
When new-crop rice first appears on the market, there may be
fluctuations of one or two \jt{俵}{tawara}; the market
consolidates (\jt{保合}{mochai}) at the floor, then rallies some
fifty or sixty tawara before pulling back two tawara or so and
entering another consolidation. At such a time, buyers grow
restless and sellers press their advantage, selling more and more.
This rice must never be sold. Those
who are short should cover their positions here. This rice is
rising from the floor. The decline is a dip within a larger
uptrend. Above all, one should consider the timing of the ceiling
price (\jt{天井直段}{tenj\=o nedan}).
\end{sokyubox}

\commentary
This chapter warns against selling into a market that is dipping
within a larger rising trend---and illustrates the point with an
example. Suppose that after the opening session (\jt{發會}{hakkai}) of
the new contract, there are minor fluctuations and the market
consolidates near the floor, then gradually firms, rallying some
fifty or sixty sen before pulling back about twenty sen and
consolidating again. In such a case, the buyers have grown weary
from the market's prolonged refusal to advance, while the sellers,
encouraged, press their sales more aggressively.
Yet in a market of this kind, one absolutely must not sell.
Anyone who has sold should cover here, for this market is in the
form of a rise from the floor---the dip is merely a decline within
the larger uptrend. If one intends to sell, one should wait for
the timing of the ceiling price; that is the essential thing.

The previous chapter warned against buying into a market that is
breaking down from a ceiling. This chapter warns the sell side
against pressing into a market that is rising from the floor. The
master's solicitous guidance is worthy of our gratitude.

% -----------------------------------------------------------------------------
%  Article 37
% -----------------------------------------------------------------------------
\article{37}{Guidance for Times of Adverse Fortune}

\begin{sokyubox}
{\small\itshape\color{sokyubar} \jp{第三十七章\quad 不利運の節心得方の事}}

\medskip
In times of adverse fortune, one must never average down on sells
or average down on buys. When one's forecast has gone wrong, close
the position at once and rest for forty or fifty days. Even after
a trade (\jt{商内}{akinai}) that has been thoroughly successful, once
the position is closed one should rest for forty or fifty days,
study the oscillation of rice prices, compare against the
Three-Position Tradition (\jt{三位の傳}{sanmi no den}), consider the
moment when the opportunity aligns, and only then re-enter the
market. No matter how great one's winning fortune, if one forgets
this rest, losses will inevitably appear when the trade is finally
closed. However, ``resting after closing a trade'' does not mean
resting absent-mindedly. One must set aside all bullish or bearish
bias and study the daily price fluctuations without letting down
one's guard. Again, when one has profited on the buy side
the previous year, one is apt to cling
stubbornly to a bullish stance---yet one must completely let go of
last year's sentiment and consider the crop conditions, the
relative abundance or scarcity of supply, and the state of market
sentiment (\jt{人氣}{ninki}) for the present year. That is the first
priority.
\end{sokyubox}

\commentary
In times of adverse fortune, one must never average down---neither
selling more to average a losing sell, nor buying more to average a
losing buy. This amounts to fighting the market and insisting on
one's own view, which violates the cardinal principle of market
success: \textit{``Follow the market.''} Therefore, when one realizes one's
forecast was wrong, one should cut the loss before it deepens and
rest for forty or fifty days---setting aside bullish or bearish
bias and, in a state of calm detachment, carefully observing the
daily fluctuations and the positions of ceiling and floor, until
one has a firm assessment of the market's condition, and only then
re-enter. Even after a thoroughly profitable trade, once the
position is closed one should rest and recuperate for forty or
fifty days, suppressing overconfidence and cultivating sound
judgment. If one forgets to rest and, intoxicated by recent
success, presses onward, one will encounter an unexpected ambush---that is, a reversal in the market---and when the position is
finally closed, it will assuredly end in loss. Yet ``resting'' here
does not mean resting idly. It means setting aside the pull of
bullish or bearish sentiment and observing the market with a Zen
mind---what might be called ``stillness within motion''
(\jp{動中の靜觀}). The details are explained in the section on
Master S\=oky\=u's ``market Zen'' and so are omitted here; but in
general, a weakness of human nature is that prior impressions
become dominant, and this weakness manifests itself with particular
frequency in the markets. To give a brief example: a person who
previously profited on the bull side tends to cling to that
experience and is inclined to buy; a person who previously profited
on the bear side has that memory impressed on the brain and is
inclined to lean bearish. But maneuvering (\jt{駈引}{kakehiki}) in
the markets is like commanding troops in battle---one must vary
one's tactics according to the moment. What matters most is to
consider the crop conditions for this year, the general economic
climate, the state of physical rice distribution, and the direction
of market sentiment.

% -----------------------------------------------------------------------------
%  Article 38
% -----------------------------------------------------------------------------
\article{38}{Guidance for Times of Profitable Fortune}

\begin{sokyubox}
{\small\itshape\color{sokyubar} \jp{第三十八章\quad 利運の節心得方の事}}

\medskip
When a trade is going well, one should first secure a reasonable
profit and take it. At that point, rest for one or two months. If
one forgets to rest, then no matter how favorable one's fortune,
losses will invariably appear when the trade is finally closed.
Flushed with victory, one whose profit is a hundred ryo begins to
think of taking two hundred; the mind drifts to a thousand, then
two thousand; lost in greed, one cannot bring oneself to cut and
losses mount. This is delusion born of avarice. In times of
adverse fortune, it is all the more so. Knowing when to cut is
essential. Be prudent and bear this in mind.
\end{sokyubox}

\commentary
This chapter carries the same meaning as the previous one, so no
separate annotation is needed. The ``one or two months'' here and
the ``forty or fifty days'' of the prior chapter amount to the same
thing. Where this chapter differs, however, is in warning against
the danger of being flushed with victory and forgetting the market's
extremes, pressing too deep into a position. One tries to profit on
both the upswing and the downswing within a short span, forcing
one's assessment and chasing the market's tail with trade after
trade---and ends up losing on both the rise and the fall. This,
in the end, is delusion born of avarice, and the chapter offers
that admonition.

% -----------------------------------------------------------------------------
%  Article 39
% -----------------------------------------------------------------------------
\article{39}{One Must Not Ride a Winning Streak}

\begin{sokyubox}
{\small\itshape\color{sokyubar} \jp{第三十九章\quad 勝に乘る可からざる事}}

\medskip
When one's forecast has been correct for several months and one has
been right eight or nine times out of ten, one must never ride the
winning streak. Simply focus on securing profits without incident.
Above all, one must not deepen one's greed and fall into delusion.
\end{sokyubox}

\commentary
Suppose one has been on the right side of the market for months---whether buying or selling---and the trade has advanced steadily,
but now the moment has arrived when the rise or the fall is
reaching its limit. If, even then, one thinks one's profits are
not yet sufficient and, out of deepening greed, holds on, the
market will reverse, and one will have merely let what could have
been a clean, uneventful profit slip away. One should make it
one's sole concern to take profit without incident. Greed is a
thing without limit.

% -----------------------------------------------------------------------------
%  Article 40
% -----------------------------------------------------------------------------
\article{40}{On Aiming for the Floor and Aiming for the Ceiling}

\begin{sokyubox}
{\small\itshape\color{sokyubar} \jp{第四十章\quad 底ねらひ天井ねらひの事}}

\medskip
Make it one's constant aim to buy at the floor (\jt{底}{soko}) and
sell at the ceiling (\jt{天井}{tenj\=o}). This should be the trader's
single-minded focus.
\end{sokyubox}

\commentary
That the essential secret of market success lies in targeting the
ceiling and the floor goes without saying, but the method of
observation is exceedingly difficult. As the various theories
would take too long to detail here, the commentary confines itself
to the master's own teachings, which hold that discerning the
ceiling and the floor requires reasoning from two factors: the
proportion of the price swing, and the number of months over which
it has unfolded.

Regarding proportions, one passage in this book states: ``In aiming
for the floor, take a hundred-tawara decline as the target; in
aiming for the ceiling, take a hundred-tawara rise as the target.''
And further:

\begin{itemize}[nosep]
\item Depending on conditions, the rise may not stop at a hundred
  tawara but extend to 130, 145, 150, or even 180 tawara.
\item In a market that rises and falls sharply, the timing of ceiling
  and floor cannot be fixed by a set number of days.
\end{itemize}

Regarding months, three doctrines exist for the turning point:
a three-month theory, a six-month theory, and a twelve-month
theory. Under the six-month theory:

\begin{itemize}[nosep]
\item In a falling market after a ceiling price, the decline proceeds
  at a rate of one to three or four tawara per month over five or
  six months.
\end{itemize}

As supporting examples:

\begin{itemize}[nosep]
\item Rice at a ceiling in the seventh or eighth month will fall
  through the twelfth month.
\item Rice consolidating at the ceiling from the first through the
  third or fourth month will certainly decline by the fifth or
  sixth month. If it has fully declined by the fifth month, expect
  a sharp rally in the sixth.
\item Rice consolidating at the floor from midwinter through the first
  and second months will certainly rise from the third and fourth
  months through the fifth and sixth. Rice at the floor in the
  seventh and eighth months should rise by the twelfth and first
  months.
\end{itemize}

Under the three-month theory:

\begin{itemize}[nosep]
\item Rice rising from the floor reaches the ceiling sharply within two
  or three months. Depending on conditions, it may decline for
  three months and rise in the fourth, or begin falling from the
  fifth month. Sometimes it falls for all six months.
\end{itemize}

The master also noted exceptions of seven or eight months and of
twelve months. Under the seven-to-eight-month theory:

\begin{itemize}[nosep]
\item Rice at a low price in the eleventh, twelfth, and first months
  should be regarded as a summer rally; it may continue rising
  through the seventh month.
\end{itemize}

Under the twelve-month theory:

\begin{itemize}[nosep]
\item In any month, when a sudden rise of ten to twelve or thirteen
  tawara occurs, regard it as a one-time ceiling. When the rise is
  five or six tawara, expect a further two months or so of
  appreciation. In a market that rises and falls sharply, the timing
  of ceiling and floor cannot be fixed.
\end{itemize}

Summarizing the master's teachings, the following principles emerge:

\textbf{Duration---normal pattern:}

\begin{enumerate}[nosep]
\item The duration is three months or more, with six months as the
  standard for a full swing between ceiling and floor.
\item Occasionally the reversal comes in the fourth month.
\end{enumerate}

\textbf{Duration---exceptions:}

\begin{enumerate}[nosep]
\item In a market that moves sharply, the ceiling or floor may form
  within one or two months.
\item Rarely, a ceiling-to-floor cycle may extend over seven or eight
  months.
\end{enumerate}

\textbf{Magnitude---normal pattern:}

\begin{enumerate}[nosep]
\item As a rule, when the swing reaches approximately one hundred
  tawara, it marks the ceiling or the floor.
\end{enumerate}

\textbf{Magnitude---exceptions:}

\begin{enumerate}[nosep]
\item When rice rises from a deep floor of forty-five or fifty
  tawara and even a minor crop shortfall occurs, the ceiling may
  not stop at a hundred tawara above but may extend to 160 or
  even 180 tawara. When rice rises from a floor of thirty-four or
  thirty-five tawara and crop conditions are average with some
  arbitrage rice (\jt{乘米}{nosei-mai}) present, a hundred-tawara
  rise should be regarded as the ceiling.
\end{enumerate}

Restating the master's teachings in modern terms and applying them
in plain language:

\begin{enumerate}[nosep]
\item A market that has been rising for three months or more, spanning
  fifty or more days, and has advanced approximately fifteen
  percent or thereabouts, is at or near the ceiling---one must not
  buy.
\item A market that has been falling for three months or more,
  spanning fifty or more days, and has declined approximately
  fifteen percent or more, is near the floor---one must not sell.
\item When a market rises or falls more than ten percent within one
  to two months, it may have reached the ceiling or the floor.
\item Observe the position and the magnitude of the ceiling and floor.
\end{enumerate}

However, the figures for duration and proportion stated above are
general principles based on the master's teachings. Readers should
adjust the proportions in light of prevailing economic conditions
and sentiment, exercising tactical flexibility (\jt{駈引}{kakehiki})
as circumstances demand. As the saying goes, the subtlety of
execution resides in the practitioner's own mind.

% -----------------------------------------------------------------------------
%  Article 41
% -----------------------------------------------------------------------------
\article{41}{How to Recognize When a Rise Has Peaked}

\begin{sokyubox}
{\small\itshape\color{sokyubar} \jp{第四十一章\quad 上げ留見様の事}}

\medskip
In a poor-harvest year, when a large swing has occurred and there
have been two margin calls (\jt{落引}{ochibiki}), one should
understand that the peak of the rise is near.
\end{sokyubox}

% -----------------------------------------------------------------------------
%  Article 42
% -----------------------------------------------------------------------------
\article{42}{When the Market Moves Neither Up nor Down and All Ten out of Ten Grow Bored}

\begin{sokyubox}
{\small\itshape\color{sokyubar} \jp{第四十二章\quad 相塲無高下十人が十人退屈の事}}

\medskip
When the market shows neither rise nor fall for two or three months,
or merely oscillates in a range-bound pattern (\jt{通相場}{kayoi s\=oba}),
all ten out of ten traders grow bored. Even the bulls begin to lean
bearish; the bears, believing their forecast is on target, press
their sells all the harder---and then the market unfailingly rises.
At that moment, both bears and bulls cry ``Just as I thought!'' and
scramble to buy back at once, driving a sharp, gapping rally. When
all ten out of ten are leaning to one side, the opposite invariably
comes. If the market moved in the direction everyone expected, it
would be an easy thing indeed---but it does not come as expected;
it is beyond the reach of calculation. This is the natural
principle of yin and yang (\jt{陰陽}{in-y\=o}).
\end{sokyubox}
